\chapter{Discussion}\label{s:discussion}

This chapter interprets the results of Chapter~\ref{s:evaluation} in light of
the research questions and the motivation set out in Chapter~\ref{s:intro}.

% ──────────────────────────────────────────────────────────────────────────────
\section{RQ1 --- Scale Shifts Failure Modes, Not Just Rates}

\subsection{The ceiling problem}

The most striking aspect of the \rqref{1} results is not the improvement in compile
and execution rates across model scales, but the ceiling on output match rate.
Even at 27B parameters, a model that compiles 95.6\% of scripts and
executes 93.2\%, only 49.6\% of generated workflows produce output that
exactly matches the reference. This means that for every two scripts the 27B model
generates and successfully executes, one produces a wrong answer.

The gap between execution rate (93.2\%) and output match rate (49.6\%) at 27B
is 43.6 percentage points, substantially larger than the same gap at 4B
(30.5 points) and 9B (35.9 points). As the model grows more capable at
producing syntactically valid BraneScript, the bottleneck shifts: rather than
failing to write code the compiler accepts, the model increasingly writes code
that is structurally correct but logically wrong. This pattern is consistent
with findings in code generation research more broadly, where larger models
outperform smaller ones on syntax-heavy benchmarks far more readily than on
semantic-correctness benchmarks requiring precise parameter passing and
multi-step reasoning.

\subsection{What the error breakdown reveals}

The error type analysis in Section~\ref{s:eval-errors} provides a finer-grained
picture of this shift. At 4B, the dominant failure category is compile-time
undefined variable errors (24.1\% of all examples), which are straightforward
structural failures: the model emits a reference to a variable it never
declared. These errors suggest the model has not reliably internalised the
BraneScript scoping rules, where every variable must be declared with
\texttt{let} before use.

At 9B, undefined variable errors persist at nearly the same rate (24.5\%),
but syntax errors drop sharply from 17.5\% to 4.8\%. The 9B model appears to
handle BraneScript grammar more reliably while still struggling with state
tracking across multi-step workflows. The rise in undefined function errors
(from 1.4\% to 3.8\%) suggests the model is becoming more ambitious in what
it attempts, calling package functions it has not verified are available,
even as its basic syntax improves.

At 27B, compile failures are almost entirely eliminated (syntax: 1.4\%,
undefined variables: 0.2\%, undefined functions: 2.4\%), but output mismatch
rises to 43.6\%. The model produces structurally valid workflows reliably and
without error, yet returns incorrect output. 

These are the hardest failures to detect and correct:
unlike a compile error, which is immediately surfaced and detectable by a machine, a
semantically wrong but syntactically valid BraneScript would silently return
wrong results to a researcher relying on the pipeline. This class of failure based on confident,
plausible-looking, but incorrect output, is arguably the
most practically dangerous failure mode for a system used in scientific data
analysis.

Considering the fact that the pipeline is aimed at domain scientists rather than 
software engineers, this shift in failure mode has practical implications for 
deployment: a system that compiles and executes successfully is not necessarily
a system that produces correct results, and may end up misleading researchers if
the output is not verified against a reference or ground truth.

\subsection{Smaller models and their practical value}
Despite the higher compilation failures seen in the 9B model compared to the 27B model,
the results hide important details: 
the smaller model is still capable of producing correct workflows in half of the cases it manages to execute.

The fact that the 9B model is substantially cheaper to run than the 27B model (as it fits
in a single Nvidia GTX 3060 GPU, a consumer-grade 5-year-old graphics card), 
suggests that for some use cases, where the cost of running a larger model is prohibitive,
a smaller model may still be a viable option, provided that additional verification 
steps are in place to correct its compilation failures.
The iterative nature of the pipeline proposed in this thesis, where the compilation step can be
verified and corrected with a feedback loop, allows for a practical trade-off between model size, cost, and accuracy.
Giving back the compilation error to the model for an iterative correction may allow the 9B model
to reach similar execution rates as the 27B model, while still being cheaper to run, and therefore 
more accessible to researchers with limited computational resources.
% ──────────────────────────────────────────────────────────────────────────────
\section{RQ2 --- Engineering Effort Dominates Model Scale}\label{s:disc-rq2}

\subsection{Prompt engineering as first investment}

The single largest effect observed in this thesis came not from model scale but
from correcting the system prompt and RAG configuration, which moved the 27B model's
execution rate from 37.9\% to 93.2\% with no change to model weights as represented in Figure~\ref{fig:baseline-first-run}.
A naive reading of the first evaluation round alone would have concluded that
even a 27B model is unreliable at BraneScript generation, and that only a
much larger commercial-grade model could close the gap. After the second round,
it is clear that this conclusion would have been wrong, and that the actual
bottleneck was grounding and instruction quality.

This has a direct implication for researchers deploying a local LLM-based
DSL translators: investing in a precise retrieval and an
unambiguous system prompt shall be the first priority, and only then model size
or further fine-tuning shall be taken in consideration.
Because direct edits to the base the prompt and to the retrieval configuration were all
revised together between the two evaluation rounds, this finding should be
read as evidence that retrieval-and-prompt engineering \emph{as a bundle}
outweighs model scale, not as an isolated measurement of RAG's individual
contribution against a no-RAG condition (Chapter~\ref{sec:threats}). It also
validates gap \gap{2}, as the missing piece for practical intent-to-WMS
translation a carefully engineered bridge between the two and not a larger model alone.

\subsection{The decomposition contribution}

The ablation in Section~\ref{s:eval-decomp} isolates one specific component
of that bridge: the intent decomposer. Removing it reduces compile rate from
66.1\% to 52.9\%, execution rate from 64.9\% to 48.1\%, and output match from
36.9\% to 26.0\% on the 9B model. The consistent degradation across all three
metrics indicates that the decomposer contributes positively at every stage of
the pipeline, not just at one particular failure mode.

A closer look at the error breakdown reveals where decomposition's benefit
is concentrated. Undefined function errors, where the model calls a function
that does not exist in the package, occur at virtually the same rate with and
without decomposition (28.5\% vs.\ 28.8\% of examples). Syntax errors, by
contrast, fall sharply from 12.3\% to 3.0\% when decomposition is enabled.
This suggests that the primary contribution of decomposition is not improved
package or function selection (that is passed to the RAG), but rather that 
breaking the intent into explicit sub-tasks gives the model a cleaner, more constrained generation
problem. Smaller, well-scoped sub-tasks appear to reduce the syntactic
complexity the model must manage in a single generation pass, resulting in
structurally more correct BraneScript. The function-level hallucinations
(calling non-existent functions within the correct package) persist at the
same rate regardless, suggesting they are driven by limitations in the model's
knowledge of package APIs rather than by input ambiguity.

It is also notable that the output match rate under no-decomposition (26.0\%)
remains well above zero, meaning the generator can still produce correct
workflows from unstructured input in roughly one in four cases. This suggests
the base model can still absorb enough BraneScript knowledge from its inputs
and produce correct outputs for single-package intents without decomposition.

Finally, a cross-model comparison further underlines decomposition's weight relative to
model scale. The 4B base model with decomposition enabled achieves 54.6\%
execution rate, while the 9B model with decomposition disabled reaches only
48.1\%. Decomposition applied to a smaller model outperforms a model one size
tier larger running without it. This implies that for deployment scenarios
where computational resources are constrained, investing in decomposition
yields more return than upgrading to a larger model.

% ──────────────────────────────────────────────────────────────────────────────
\section{RQ3 --- Fine-Tuning Degrades Performance Under Current Conditions}

Both fine-tuned variants collapse to near-zero compile and execution rates on
the held-out test set, far below the 9B base model they were derived from. The
comparison between ep5 (five epochs) and ep3-new (three epochs) rules out
training duration as the cause: the results are nearly identical, confirming
this is a systematic failure rather than an overfitting artefact.

\subsection{How SFT training is structured}

Each training example is formatted as a three-part conversation. The
\emph{system message} contains the BraneScript instructions and constraints
given to the model. The \emph{user message} contains the natural-language
intent and any retrieved package context. The \emph{assistant message}
contains the reference BraneScript code to be generated. These three parts
are concatenated into a single token sequence and fed to the model.

During training, the model is asked to predict each next token given all
previous ones --- a standard language modelling objective. However, only
the \emph{assistant} tokens contribute to the loss: the system and user
tokens are masked and receive zero gradient. This is called response-only
masking, and its purpose is to teach the model exclusively to generate
correct BraneScript, not to reproduce the instructions or the user's request.
The loss function used is cross-entropy: at each assistant token position,
the model outputs a probability distribution over its vocabulary, and the
loss penalises assigning low probability to the correct next token. The
LoRA adapter weights are updated via backpropagation to reduce this loss.

\subsection{Why memorisation occurred}

Three compounding factors drove the training into memorisation rather than
generalisation.

\paragraph{Dataset too small relative to model capacity.}
Qwen3.5-9B is a model pre-trained on trillions of tokens. Its LoRA adapter,
even with a conservative rank of 16, still contains millions of trainable
parameters. With only 497 training examples, there is vastly more model
capacity than data. Under these conditions gradient descent finds the path
of least resistance: memorising the 497 specific input--output mappings
exactly, rather than abstracting reusable patterns. Studies on fine-tuning
small and mid-sized LLMs find that measurable generalisation requires at least
1{,}000--5{,}000 diverse, high-quality examples for code generation
tasks~\cite{guo2025unveiling,chen2025massive}; 497 falls below the threshold
at which any generalisation benefit over the base model can be expected.

\paragraph{Learning rate too aggressive.}
The SFT learning rate of $2 \times 10^{-4}$ is high for a model of this size.
A high learning rate causes large weight updates per step, snapping the adapter
tightly to the training examples in very few steps. The training loss collapsed
to near zero within 0.3 epochs --- after seeing roughly 150 examples for the
first time --- confirming that memorisation was complete almost immediately.
A lower learning rate (e.g.\ $1 \times 10^{-5}$) would have slowed this
process, potentially allowing the model to build more generalisable
associations before the adapter overfit.

\paragraph{No mechanism to stop memorisation.}
There was no early stopping based on validation loss, no weight decay, and
no dropout on the adapter. Nothing penalised or interrupted memorisation once
it began. The trainer continued updating weights for the full number of epochs
even though loss had already reached zero, reinforcing the memorised mappings
further.

\subsection{Why memorisation causes test-set failure}

The LoRA adapter encodes 497 specific context-window patterns: for each
training example, it learned ``when the system prompt and this specific user
intent are in context, output this specific BraneScript''. On the test set,
the system prompt is the same but the user intents are new --- ones the model
has never seen. The adapter has no memorised answer for them.

Crucially, the adapter does not simply stay inactive on unseen inputs: its
weights are always applied on top of the base model, and they actively push
the generation away from what the base model would naturally produce. The
result is a generation that is neither the clean output of the pre-trained
base model nor a valid BraneScript: the adapter disrupts the base model's
competent zero-shot behaviour without replacing it with anything useful.
This explains why the fine-tuned models (2--4\% compile rate) perform far
worse than the base model they were derived from (49\% compile rate).

\subsection{Symptoms of memorisation in the generated outputs}

The character of the generated outputs confirms this picture. Three patterns
appear consistently across both fine-tuned variants.

\paragraph{Verbose inline reasoning.}
Training examples contain only intent--BraneScript pairs with no comments or
intermediate reasoning; the average training output is 352~characters.
Generated outputs average 1{,}800--2{,}000~characters --- roughly five times
longer. On unseen intents the adapter provides no memorised continuation, and
the base model's pre-trained code generation behaviour floods back: multi-line
comment blocks, repeated \texttt{commit\_result} calls with conflicting
arguments, and self-contradictory remarks about the correct approach.

\paragraph{Context contamination.}
Although the system and user tokens receive zero loss during training, the
model still sees them in the context window for every training example. With
loss collapsing immediately, the model effectively memorises each
training example as a complete context-window pattern, conflating prompt
content with assistant output. On unseen inputs it reproduces verbatim
fragments of the system prompt as inline comments mid-script --- phrases such
as \texttt{\# Do NOT attempt to re-implement the masking logic} or
\texttt{\# package function outputs}, drawn directly from the system prompt's
constraint list, appear inside the generated BraneScript. This is a
consequence of full-context memorisation rather than a failure of the
masking scheme itself.

\paragraph{Truncated scripts and EOF errors.}
95--98\% of compile failures are \texttt{reached end-of-file unexpectedly}
errors: scripts are syntactically plausible up to the point where they are
cut off mid-block. The model has a strong prior for starting a BraneScript
structure (learned from seen examples) but no reliable prior for concluding
one on a novel intent, so generation simply runs out of coherent content
before closing all open blocks.

\subsection{Implications and scope of the finding}

The near-zero compile rates should be read as a consequence of insufficient
training scale rather than as evidence that BraneScript is unlearnable via
SFT. The most direct remedy is a larger, more diverse training set on the
order of several thousand examples covering the full range of package
combinations and intent types. Complementary measures --- a lower learning
rate, early stopping on validation loss, and data augmentation through intent
paraphrasing --- would further reduce the risk of memorisation even with a
modestly sized dataset. Additionally, fine-tuning was limited to the 9B model
due to GPU allocation constraints on Snellius; larger models (27B, 70B) with
stronger zero-shot instruction-following capabilities may require fewer
examples to generalise, making them potentially more robust to the data
scarcity observed here. All of these remain open directions for future work.

% ──────────────────────────────────────────────────────────────────────────────
\section{RQ4 --- Semantic Caching as a Practical Reproducibility Layer}\label{s:disc-rq4}

The evaluation of the semantic cache yields two sets of findings that must be
read together. On one hand, the cache achieves a 99.1\% hit rate and a
98.3\% correct rate at the 0.92 threshold over the paraphrase test set,
indicating that sentence-embedding similarity is an effective proxy for
intent equivalence. On the other hand, the threshold sweep reveals an irreducible
false-positive floor of approximately 0.9\% for the same threshold value, attributable to the inability
of general-purpose embeddings to distinguish intents that share the same
topic but differ in specific parameter values.

The first finding addresses the core design goal of the cache: to return the
same verified workflow for semantically equivalent intents without repeating
the generation step ensuring reproducibility. The observed hit and correct rates confirm that this
goal is met under the evaluation conditions. The second finding establishes
a boundary on when the cache should be trusted: its error rate is hardly reducible 
to zero by threshold tuning alone. The two findings together suggest that the cache
is best understood as a \emph{performance optimisation with a bounded, known error rate} rather
than a reproducibility guarantee in the strict scientific sense.

\subsection{The false-positive floor}

The cache threshold sweep reveals a
false-positive rate that remains stable at approximately 0.9\% across
thresholds from 0.80 to 0.92, and only drops to 0.1\% at a threshold of
0.98, at which point the hit rate collapses from 99.1\% to 59.1\%.
This behaviour indicates that the false positives are not a function of the
threshold being too low, but of the embedding model mixing structurally
similar intents regardless of parameter values. Unlike the \rqref{1}/\rqref{2}
findings above, this is not attributable to fixable engineering choices: a
general-purpose sentence embedding model encodes what an intent is
\emph{about}, not the exact parameter values within it. Two intents that
differ only in a threshold value, a column name, or a dataset identifier are,
for most practical purposes, requests for different computations, yet they
are nearly indistinguishable in embedding space. Threshold tuning is therefore
not a completely viable path towards the elimination of false positives.

This is a direct, empirical instance of gap \gap{3}, where no agreed definition of intent
reproducibility exists (Chapter~\ref{s:background}): the semantic cache implicitly
assumes that ``reproducible'' means ``the same workflow for the same
\emph{propositional} intent'', while for scientific correctness it must mean
``the same workflow for the same intent \emph{including all parameter values}''.

A semantic cache that weights parameter values more heavily than topic similarity 
would be a more appropriate solution for scientific reproducibility, but it would require a
different embedding model or a custom similarity function that is aware of the DSL's 
parameter semantics. This is a potential avenue left for future work.

\subsection{The operating point and its practical meaning}

The remaining 0.8\% false-positive rate from the 0.92 threshold means that for every 125 cache hits, approximately
one will return a workflow intended for a semantically similar but distinct
intent. In a scientific computing context, this is not a negligible risk, as a
wrong workflow executed silently on a dataset could produce results that look
plausible but are derived from the wrong analysis. For this reason, the cache should
not be used as the sole quality gate.

It is worth noting, however, that the threshold can be treated as a
confidence dial rather than a fixed parameter. Raising the threshold does not
simply reduce false positives by a small margin, but it causes the cache to
abstain on lower-confidence matches entirely, routing them to a fresh
generation call instead. At 0.96, for example, the false-positive rate drops
to 0.65\% while the hit rate falls to 88.1\%, meaning that roughly 12\% of
requests that would have been served from cache are instead regenerated. For
deployments where correctness is critical and generation cost is acceptable,
a higher threshold is therefore a principled strategy: uncertain matches will
receive a new and independently generated workflow rather than a cached one of
unknown validity. The threshold, in this sense, encodes a policy choice about
the acceptable trade-off between latency and confidence, and the sweep data
in Section~\ref{s:eval-cache} provides the empirical basis for making that
choice explicitly.


% ──────────────────────────────────────────────────────────────────────────────
\section{Revisiting the Choice of Brane}

The motivation for using Brane as the testbed for this work
(Chapter~\ref{s:background}) rested on two properties: its support for
federated workflows in which the party executing a computation need not own
the underlying data, and its enforcement, via typed, sandboxed packages
of that separation at the platform level rather than by convention. 

The results in this thesis evaluate translation accuracy and reproducibility, 
but the selection of an underlying architecture that makes evaluation possible under private conditions
is itself evidence for a privacy-centered design. 
The translation problem studied here is therefore additive to Brane's existing privacy guarantees
rather than a replacement for them: a correct BraneScript workflow inherits
Brane's data-locality enforcement automatically, simply by virtue of being
expressed in BraneScript rather than in a Python script with unrestricted file
access.

The static typing of BraneScript has a secondary benefit that emerged
empirically through this thesis: it makes evaluation cleaner than it would be
for a dynamically typed target language. Because the compiler rejects undefined
variables, undefined functions, and type mismatches at compile time, compile
rate is a meaningful and automatically computable proxy for a large class of
semantic errors, without requiring execution. This is not a property of most
target languages used in comparable NL-to-code work (Python, SQL, Bash), where
many structural errors only surface at runtime or not at all. The strictness
that makes BraneScript difficult to generate correctly is the same property
that makes incorrect generation immediately detectable.

% ──────────────────────────────────────────────────────────────────────────────
\section{Summary}

Taken together, the results suggest a layered picture of where accuracy comes
from in intent-based scientific workflow generation.

Retrieval and prompt engineering account for the largest single improvement
observed, and should be the first investment for anyone building a similar
system for a low-resource DSL. The decomposer contributes a further, consistent
improvement by converting vague multi-clause intents into structured retrieval
queries.

Scaling the base model from 4B to 27B brings substantial compile and execution
rate improvements, but the output match ceiling at 49.6\% reveals that scale
alone cannot solve the semantic accuracy problem. The dominant failure mode at
27B is no longer structural but semantic: the model generates plausible,
executable workflows that return the wrong answer, a failure class that is
harder to detect and more dangerous in practice than a compile error.

Semantic caching provides a strong and computationally cheap reproducibility
mechanism for repeated or paraphrased intents, but has an irreducible
false-positive floor imposed by the limitations of sentence embeddings as a
proxy for computational equivalence. The threshold of 0.92 represents an
acceptable operating point and not a principled solution to the parameter
sensitivity problem. System administrator may choose to tune the threshold to
a higher value if correctness is critical and generation cost is acceptable.
